本节选定 Grothendieck 宇宙 $\mathcal{U}$ 和交换环 $\Bbbk$. 相对于 $\mathcal{U}$, 所有的环都默认是小的. 对于 $\Bbbk$-线性范畴之间的函子同样默认 $\Bbbk$-线性, 如定义 \ref{def:cat-k-linear}, \ref{def:k-linear-cat} 所述. 当 $\Bbbk = \Z$ 时, $\Bbbk$-线性也就是加性.

对于 $\Bbbk$-线性范畴 $\mathcal{C}$ 和 $\mathcal{D}$, 符号 $\mathrm{Fct}(\mathcal{C}, \mathcal{D})$ 或 $\mathcal{D}^{\mathcal{C}}$ 代表所有函子 $F: \mathcal{C} \to \mathcal{D}$ 构成的范畴. 尽管 $\mathrm{Fct}(\mathcal{C}, \mathcal{D})$ 未必是 $\mathcal{U}$-范畴, 由于我们不在其中作极限等构造, 这不会对往后的讨论造成任何困难.

\begin{definition}
	\index[sym1]{Fctc@$\mathrm{Fct}^c, \mathrm{Fct}_c$}
	设 $\mathcal{C}$ 和 $\mathcal{D}$ 为 $\Bbbk$-线性范畴.
	\begin{itemize}
		\item 设 $\mathcal{C}$ 和 $\mathcal{D}$ 余完备, 定义 $\mathrm{Fct}(\mathcal{C}, \mathcal{D})$ 的全子范畴 $\mathrm{Fct}^c(\mathcal{C}, \mathcal{D})$, 由全体保小 $\varinjlim$ 的函子组成;
		\item 设 $\mathcal{C}$ 和 $\mathcal{D}$ 完备, 定义 $\mathrm{Fct}(\mathcal{C}, \mathcal{D})$ 的全子范畴 $\cate{Fct}_c(\mathcal{C}, \mathcal{D})$, 由全体保小 $\varprojlim$ 的函子组成.
	\end{itemize}
	一些文献称 $\mathrm{Fct}_c$ (或 $\mathrm{Fct}^c$) 的对象为连续 (或余连续) 函子.
\end{definition}

我们沿用 \S\ref{sec:otimesL} 的符号, 对任意 $\Bbbk$-代数 $A$ 和 $B$, 定义
\begin{align*}
	(A, B)\dcate{Mod} & := (A, B)\text{-双模范畴}, \\
	A\dcate{Mod} & := \text{左}\; A\text{-模范畴} \simeq (A, \Bbbk)\dcate{Mod}, \\
	\cated{Mod}B & := \text{右}\; B\text{-模范畴} \simeq (\Bbbk, B)\dcate{Mod}.
\end{align*}
它们是 $\Bbbk$-线性范畴的简单例子.

回忆以下的模论常识: $(A, B)$-双模 $P$ 典范地诱导保小 $\varinjlim$ 的函子
\[ P \dotimes{B} (\cdot): B\dcate{Mod} \to A\dcate{Mod}, \quad (\cdot) \dotimes{A} P : \cated{Mod}A \to \cated{Mod}B, \]
而 $(B, A)$-双模 $Q$ 典范地诱导保小 $\varprojlim$ 的函子
\[ \Hom_{\cated{Mod}B}(Q, \cdot): B\dcate{Mod} \to A\dcate{Mod}, \quad \Hom_{\cated{Mod}A}(Q, \cdot): \cated{Mod}A \to \cated{Mod}B, \]
若改取函子 $\Hom(\cdot, R)$, 其中 $R$ 是 $(B, A)$-双模, 则其定义域须换为相反范畴如 $(B\dcate{Mod})^{\opp}$ 等, 使函子依然保小 $\varprojlim$.
